% Modernised reading — fonds Alexandre Grothendieck, shelfmark 119, pages 141-313
% (the transcribed batches only: 8, 13, 14 and 16).
%
% This is an INTERPRETATION, not a transcription. Everything the critical
% apparatus carried moves into footnotes. In any dispute of substance the
% transcription governs: whoever cites Grothendieck must cite batch-NN.fr.tex.
%
% Pass: Opus 5.5 (claude-opus-5-5), 2026-10-03 — first pass, unchecked against
% the pages by a human. Held to the standard calibrated on folder 115 (Opus 5),
% but not measured against that reading.
%
% Scope. Folder 119 has 313 pages in sixteen batches; four are transcribed
% (8: pp. 141-160; 13: pp. 241-260; 14: pp. 261-280; 16: pp. 301-313). The
% Esquisse d'un programme itself (pp. 2-59 and 78-135, two copies) is printed
% (Schneps & Lochak, 1997) and is neither transcribed nor restated here. The
% other untranscribed items are listed in issue #45 and named only in the spine
% section. Third-party pieces (Orsay commission papers, CNRS letter of 1972,
% Serre's résumé) are mentioned no further than the transcription's notes.
%
% Notation fixed for the whole document. Tame ("modéré") spaces are compact
% unless said otherwise; I = [0,1]. B is a "bordant" subspace, (C, Ċ) a collar
% or a cone, Y a closed tame subspace, (V, V̇) a tubular neighbourhood. For the
% combinatorial pages: S_n (fraktur) symmetric groups, ω(S) the orientation
% torsor of a finite set, ∧ the contracted product of Z/2-torsors; his GP(1,k)
% is written PGL_2(k), his Aff(n,k) is AGL_n(k).
%
% Substantive departures from the pages, each footnoted where it occurs:
% — p. 147, conjecture 1 c): read literally (homeomorphisms of B x I fixing
%   B x 0), it is a statement about a pseudo-isotopy space, which is not
%   contractible in general (Hatcher-Wagoner); the reading says so.
% — p. 148: his « (Alexandroff ?) » is the Alexander trick (J. W. Alexander);
%   conjecture II c) is true by that trick, and the reading says so.
% — p. 150: X \ (V̇ \ V) on the page is read X \ (V \ V̇).
% — p. 155: « sous-groupe d'ordre 10 » of a group of order 72 is impossible as
%   read; not reconstructed.
% — p. 158: his « c implique card V = 9 » also needs a).
% — p. 304: « [16, 17] » for the programme of algebraic geometry is probably
%   [15, 16]; p. 306: the pages of [13] repeat those of [12].
%
% What the folder announces and never establishes: the « Th » of p. 145 has no
% proof; the eleven "∞-isotopic equivalences" of p. 149 and the conjectures of
% pp. 147-152 are a programme; the plan of pp. 154-160 is a plan. Part B of the
% 1968 text is announced in ink on the first state and only written in the
% second.
\documentclass[11pt,a4paper]{article}
\input{../preamble/grothendieck}

\folder{119}
\pages{141}{313}
\dating{1968-1991}
\watermark{Édition de démonstration\\interprétation personnelle de l'œuvre}
\foldertitle{Esquisse d’un programme [dossier constitué en vue d'une candidature au CNRS (1984)] : copies de tapuscrits et de tapuscrits annotés (1972-1974, 1978-1979, 1984, 1990-1991, s.d.), tirés à part (1971), notes et copies de note manuscrites (1986, s.d.), lettres (1968, 1972, 1991). — lecture modernisée du dossier entier}

\begin{document}

\section*{Résumé}

\begin{resume}
Ce dossier est celui qu'il a constitué en 1984 pour entrer au CNRS. On
s'attendrait à n'y trouver que l'\emph{Esquisse d'un programme}, le texte
célèbre qu'il écrit alors pour présenter ce qu'il compte faire. Elle y est,
en deux exemplaires, mais elle est imprimée depuis 1997 et cette édition ne la
reprend pas. Le reste du dossier est autre chose : une trentaine d'années de
sa vie déposées dans une chemise. Cette lecture ne couvre que la partie de ce
reste qui a été transcrite — quatre lots sur seize, environ soixante-dix
pages du dossier, dont une cinquantaine de sa main ou de sa machine — et elle
le dit à chaque fois que cela
compte.

Il y a d'abord trois notes mathématiques manuscrites, postérieures à
l'\emph{Esquisse} (1984--1986). La première se demande ce que devient un
espace quand on le « quotiente » par un feuilletage — quand on le découpe en
feuilles, comme un livre en pages, et qu'on veut parler de l'ensemble des
feuilles. Cet ensemble est en général très mauvais comme espace : une feuille
peut s'enrouler indéfiniment et passer aussi près qu'on veut de toutes les
autres. Il propose de le remplacer par un \emph{topos}, objet qui retient
l'information perdue par le quotient naïf, et énonce une équivalence entre
les deux façons de décrire la situation. C'est aujourd'hui le topos classifiant
du groupoïde d'holonomie.

La deuxième, un « Addendum » à l'\emph{Esquisse}, revient sur la
\emph{topologie modérée} : une géométrie où l'on interdirait les monstres de la
topologie générale — courbes qui remplissent un carré, spirales infinies —
pour ne garder que des figures qui se laissent découper en morceaux simples.
Il y formule des conjectures sur les \emph{colliers} (une bande $B \times
[0,1]$ le long d'un bord), les \emph{cônes} (le voisinage d'un point
ressemble au cône sur un petit espace, comme le voisinage du sommet d'un
cornet ressemble au cône sur un cercle) et les \emph{voisinages tubulaires}.
Le principe de toutes est le même : les choix qu'on fait pour construire ces
voisinages devraient être « anodins », c'est-à-dire que l'espace de tous les
choix possibles devrait être contractile, de sorte que rien de ce qu'on en
déduit ne dépende du choix. Une partie de ce qu'il conjecture est devenue un
théorème dans la théorie des structures o-minimales ; une autre, lue à la
lettre, se heurte à un résultat connu, et la lecture dit lequel.

La troisième est le plan d'un texte de combinatoire sur les ensembles à au
plus six éléments : les groupes de permutations de ces petits ensembles sont
« accidentellement » isomorphes à des groupes de matrices sur de petits corps
finis, et le groupe des permutations de six objets a une symétrie cachée qui
échange les objets avec d'autres objets, au nombre de six aussi. Il veut en
faire une géométrie, avec icosaèdres et plans affines sur le corps à trois
éléments.

Le dossier contient ensuite un texte de juin 1968, « Le maître-enseignant et
le maître-chercheur dans l'Université d'aujourd'hui et de demain », avec la
lettre par laquelle il le propose au journal \emph{Le Monde}. C'est un texte
politique, écrit au milieu du mouvement de mai, et il ne demande aucune
mathématique. Sa thèse est simple et a gardé son tranchant : enseigner et
chercher sont deux métiers, apparentés mais distincts ; l'Université les
confond, ne recrute que sur la recherche, et méprise ceux qui ne font
« que » enseigner, au détriment des étudiants d'abord. Il propose de séparer
les deux carrières, de les mettre à égalité, et — c'est le seul point qu'il
tient pour essentiel — de donner aux étudiants une part réelle dans les
décisions qui les concernent. Le chercheur qu'il décrit, et qui parle de
mathématiques au premier congénère venu, « dans un métro » s'il le faut, est
un autoportrait à peine voilé.

Enfin, un curriculum vitae d'environ 1971, qui va de sa naissance à Berlin
et des camps d'internement jusqu'à son départ de l'IHES et à la fondation de
\emph{Survivre}, et qui se termine sur une phrase qu'on ne s'attend pas à
lire dans un dossier de candidature : la seule valeur de son œuvre de
mathématicien, dit-il, est de donner du poids à son témoignage pour la survie
de l'espèce.

Ce qui relie ces pièces n'est pas un sujet mais une attitude. L'addendum
cherche des énoncés où les choix soient « anodins », le texte de 1968 des
structures qui « s'ajusteront d'elles-mêmes » une fois le bon principe acquis.
Dans les deux cas, il se méfie de ce qu'on fixe par décret et cherche ce qui
tient tout seul.

\keywords{foliation, holonomy groupoid, classifying topos, open geometric morphism, tame topology, o-minimal structure, collar theorem, pseudo-isotopy, Alexander trick, local conic structure, mapping cylinder neighbourhood, Thom–Mather stratification, intrinsic stratification, exceptional isomorphisms of symmetric groups, outer automorphism of S6, Steiner triple system, affine plane over F3, May 1968 university reform}
\end{resume}

\pagerange{141}{143}
\section*{Le dossier, ce que cette lecture en couvre, et ses notations}

Le dossier 119 compte 313 pages. Les deux exemplaires de l'\emph{Esquisse
d'un programme} (janvier 1984) en occupent cent seize (pages 2 à 59 et 78
à 135) ; ce texte est publié\footnote{L.~Schneps et P.~Lochak (éd.),
\emph{Geometric Galois Actions~1}, London Math. Soc. Lecture Note Series~242,
Cambridge, 1997. Il n'est ni transcrit ni reformulé ici.} et cette lecture
n'en dit rien de plus. Le reste a été relevé page par page sur le fac-similé,
et seule une partie en a été transcrite. Cette lecture porte sur ces parties
transcrites, et sur elles seules :

\begin{itemize}
\item pages 141 à 160 : trois notes mathématiques manuscrites
  (1984--1986), précédées d'un rapport d'activité déjà publié ailleurs ;
\item pages 241 à 280 : le texte de juin 1968 sur le maître-enseignant et le
  maître-chercheur, en deux états, la lettre à Hubert Beuve-Méry, et des
  pièces collectives de la commission d'Orsay ;
\item pages 301 à 313 : son curriculum vitae, en deux exemplaires.
\end{itemize}

Ce qui n'est pas transcrit — des tirés à part de N.~Saavedra Rivano (1971), un
second rapport d'activité (décembre 1984), le polycopié \emph{Fonctions
holomorphes} (1973--1974), le programme de C4 de 1978--1979, deux exemplaires
de l'\emph{Esquisse thématique} (1972), une lettre de Ronald Brown (1991) et un
résumé de cours de J.-P.~Serre — n'est pas lu ici. Une lecture faite sans ces
pages ne peut pas dire ce qu'elles changeraient, et ne le prétend pas.

Les pages 141 et 142 portent un « Rapport d'activité » dactylographié, non
daté, titré de sa main ; il est déjà transcrit par ailleurs et la
transcription ne le reprend pas.\footnote{Transcription en ligne du Cercle
Grothendieck / CSG, cote U18d. La note de la transcription le résume ainsi :
son enseignement depuis 1978 (C4, DEA, options de DEUG), conçu comme une
initiation à la recherche où la difficulté principale des étudiants est la
\emph{verbalisation} ; ses directions de thèses (topologie des surfaces) ;
ses sujets de réflexion — géométrie algébrique anabélienne et tour de
Teichmüller, topologie modérée et dévissage des structures stratifiées,
fondements de l'algèbre homotopique.} On le mentionne parce qu'il nomme
exactement les sujets des notes qui suivent : la topologie modérée et le
dévissage des structures stratifiées, c'est l'« Addendum » des pages 147
à 152. La page 143 est une chemise, d'une autre main sans doute : « Notes de
math. oubliées par Jean Malgoire », qui couvre les pages 144 à 160.

Le dossier n'a donc pas de fil mathématique continu, et cette lecture n'en
invente pas. Elle a cinq stations, chacune autonome :

\begin{enumerate}
\item le topos quotient d'un feuilletage (pages 144--145, mars 1986) ;
\item l'« Addendum » à l'\emph{Esquisse} : colliers, cônes, voisinages
  tubulaires en topologie modérée (pages 147--152) ;
\item les « Digressions combinatoires » : géométrie des ensembles à au plus
  six éléments (pages 154--160) ;
\item le texte de 1968 et la lettre au \emph{Monde} (pages 241--280) ;
\item le curriculum vitae (pages 301--313).
\end{enumerate}

Pour les trois premières, qui sont des mathématiques, la lecture énonce ce qui
est vrai et dit en note ce que porte la page. Pour les deux dernières, qui ne
le sont pas, elle est une présentation fidèle en français d'aujourd'hui, avec
le contexte que le texte suppose connu ; elle n'y corrige rien, sauf des
références bibliographiques.

\emph{Notations, valables pour les trois stations mathématiques.} Un espace
\emph{modéré} est un espace de la topologie modérée de l'\emph{Esquisse} ; la
lecture le prend au sens, précisé depuis, d'un ensemble définissable dans une
structure o-minimale (par exemple semi-algébrique), compact sauf mention
contraire.\footnote{Les pages ne définissent pas « modéré » : elles
supposent connue la section de l'\emph{Esquisse} qui en esquisse la notion.
Le choix du cadre o-minimal est le nôtre ; il est fait parce que c'est là que
les énoncés des pages 147 à 152 ont été depuis démontrés ou réfutés.} On
note $I = [0,1]$. Pour les pages combinatoires, $\mathfrak{S}_n$ est le groupe
symétrique, $\omega(S)$ l'ensemble à deux éléments des orientations d'un
ensemble fini $S$, $\wedge$ le produit contracté des ensembles à deux
éléments ; ce qu'il note $\mathrm{GP}(1,k)$ est écrit
$\mathrm{PGL}_2(k)$, et son $\mathrm{Aff}(n,k)$ est écrit
$\mathrm{AGL}_n(k)$.

\pagerange{144}{145}
\section*{Le topos quotient d'un feuilletage (pages 144 et 145, mars 1986)}

La note est datée « Mars 1986 » et titrée en marge « Feuilletages », avec
une question à lui-même : « Mettre où ? ».

\subsection*{Le problème}

Soit $E$ un espace muni d'une \emph{relation d'équivalence locale} $Q$ :
l'exemple type est une variété feuilletée, où deux points sont localement
équivalents quand ils sont sur la même plaque d'une carte feuilletée.
L'espace des feuilles $E/Q$, quotient topologique naïf, est en général
dégénéré — pour le feuilletage de Kronecker du tore par des droites de pente
irrationnelle, chaque feuille est dense et le quotient porte la topologie
grossière. Il propose de remplacer ce quotient par un \emph{topos quotient},
et note au passage que $E$ pourrait déjà être un topos.

Il donne un cas modèle : les quotients obtenus sont des « espaces » de la
forme $(X, G)$, où $X$ est un espace et $G$ un groupe discret opérant sur $X$,
pas nécessairement librement, $G$ étant le groupe fondamental d'une feuille.
C'est la situation d'une \emph{suspension} : si toutes les feuilles sont
homéomorphes à un même $F$ et qu'une application $E \to F$ induit des
homéomorphismes de chaque feuille sur $F$,\footnote{Une partie de la
condition est illisible sur la page (« et \ldots{} : $E \to F$ induisant des
iso $F_i \to F$ »). La lecture la complète par l'hypothèse qui rend le cas
modèle correct ; ce complément est le nôtre.} le feuilletage est localement
un produit transverse, l'holonomie est une action de $G = \pi_1(F)$ sur une
transversale $X$, et le topos quotient attendu est le topos $B_G(X)$ des
faisceaux $G$-équivariants sur $X$ — le « quotient homotopique » $X /\!/ G$,
qui garde trace des stabilisateurs là où l'action n'est pas libre. Une note
marginale presque illisible porte le mot « holonomie ».\footnote{Lu
« \emph{Notion} \ldots{} de l'"holonomie" », le premier mot douteux.}

\subsection*{L'énoncé}

Sous le sigle « Th », il énonce ceci. Soit $C$ la catégorie des couples
$(X, R)$, où $X$ est un espace topologique et $R$ une relation d'équivalence
locale sur $X$, localement ouverte, à fibres localement connexes, les flèches
étant les applications continues compatibles aux relations. Soit $C'$ la
catégorie des triplets $(X, \mathcal{Y}, f)$, où $X$ est un espace
topologique, $\mathcal{Y}$ une « multiplicité topologique », et
\[
  f \colon \mathrm{Top}(X) \longrightarrow \mathcal{Y}
\]
un morphisme de topos, issu du topos des faisceaux sur $X$, localement ouvert
et à fibres localement connexes. Alors
\[
  C \;\simeq\; C' .
\]
Il précise que $C'$ est a priori une 2-catégorie mais que ses catégories de
morphismes $\mathrm{Hom}(\xi, \eta)$ sont \emph{discrètes}, de sorte qu'il
s'agit bien d'une équivalence de catégories ordinaires.\footnote{Plusieurs
mots de l'énoncé sont biffés ou illisibles, en particulier une troisième
condition sur $R$ et sur $f$ (« à fibres loc. connexes et \ldots »). Les
lectures « \emph{mult. top.} » et « \emph{flèches} » sont incertaines. La
lecture ne restitue pas la condition manquante.}

Ce qu'il faut entendre est que le foncteur $C \to C'$ envoie $(X, R)$ sur le
morphisme quotient de $\mathrm{Top}(X)$ vers le topos quotient de $X$ par $R$,
et que réciproquement $R$ se lit sur $f$ comme la relation « être dans la
même composante connexe locale d'une fibre ». La page ne donne aucune
démonstration, et la lecture ne prétend pas en avoir une : c'est un énoncé
annoncé.\footnote{Sur les mots. « Multiplicité » est le terme qu'il emploie
dans ces années pour des objets de type champ ou topos généralisant les
orbifolds. Les notions de morphisme de topos \emph{ouvert} et
\emph{localement connexe} sont dans le mémoire de Joyal et Tierney,
\emph{An extension of the Galois theory of Grothendieck} (Memoirs AMS 309,
1984) ; la page ne cite rien. Le groupoïde d'holonomie (Ehresmann ;
Winkelnkemper, 1983) et l'usage de groupoïdes et de classifiants pour les
feuilletages (Haefliger, depuis les années 1950 ; \emph{Groupoïdes d'holonomie
et classifiants}, Astérisque 116, 1984) précèdent la note ; la description du
« topos classifiant d'un feuilletage » comme topos classifiant du groupoïde
d'holonomie, et l'étude systématique des topos de groupoïdes étales, sont
dues pour l'essentiel à Moerdijk et ses collaborateurs à partir de la fin
des années 1980. Rien sur la page ne dit qu'il connaissait ces travaux.}

Il ajoute des variantes pour les variétés feuilletées de classe $C^{r}$,
$r \in [0, +\infty] \cup \{\omega\} \cup \{\omega_{\mathbb{C}}\}$ —
c'est-à-dire du topologique à l'analytique complexe — et une question :
\emph{quelles variétés} (au sens de ces multiplicités) obtient-on comme
quotients de feuilletages ? La page la laisse ouverte, et une marge
illisible semble y revenir (« \ldots{} cette lacune \ldots{} feuilletages »).

\pagerange{147}{152}
\section*{Addendum à l'\emph{Esquisse} : colliers, cônes, voisinages tubulaires (pages 147 à 152)}

Six feuillets numérotés 1 à 6 de sa main, non datés. Le titre est
« Addendum Esquisse d'un Programme », et la première ligne en annonce
l'objet : le \emph{voisinage tubulaire} d'une partie modérée d'un espace
modéré compact, avec pour cas évident celui d'un point.\footnote{Le mot
« modérée » remplace « quelconque », biffé : la restriction aux parties
modérées est voulue, et elle est essentielle — pour une partie quelconque
d'un espace, même agréable, il n'y a pas de voisinage tubulaire.}

L'idée directrice, qu'il répète à chaque paragraphe, est que les choix faits
pour construire un voisinage — un collier, un cône, une paramétrisation —
doivent être \emph{anodins} : l'espace de ces choix doit être contractile.
C'est la forme précise de « ne dépend pas du choix » quand ce qu'on construit
n'est pas une classe mais un objet, et que l'on veut pouvoir en transporter
les automorphismes.

\pagerange{147}{147}
\subsection*{I. Colliers}

Soit $B$ une partie modérée « bordante » d'un espace modéré $X$ — la page ne
définit pas le mot, qui doit vouloir dire que $B$ est placée dans $X$ comme un
bord.\footnote{Au-dessus de « bordant », il a écrit puis biffé « bon
"bord" ». La conjecture 1 a) ci-dessous, qui affirme l'existence d'un collier
pour toute partie bordante, montre que « bordant » est une notion intrinsèque,
distincte de l'existence d'un collier ; la page ne dit pas laquelle.}
Un \emph{collier} autour de $B$ est un voisinage $C$ de $B$ dans $X$ avec un
homéomorphisme modéré
\[
  \varphi \colon C \xrightarrow{\ \sim\ } B \times I, \qquad
  \varphi(b) = (b, 0) \ \ (b \in B),
\]
et l'on pose $\dot C = \varphi^{-1}(B \times \{1\})$.\footnote{La page porte
$\varphi(B \times \{1\})$ ; l'exposant $-1$ n'y est pas visible, mais il est
demandé par le sens, $\dot C$ étant une partie de $C$.} Le collier est
\emph{strict} si $\dot C$ est à son tour bordante dans
$X \smallsetminus (C \smallsetminus \dot C)$ ; il est \emph{rigidifié}
quand on retient $\varphi$ et pas seulement le couple $(C, \dot C)$.

\textbf{Conjecture 1.} Soit $B$ bordante dans $X$.
\begin{enumerate}
\item[a)] Il existe un collier, donc aussi un collier strict.
\item[b)] L'espace des colliers stricts est contractile.
\item[c)] L'espace des rigidifications d'un collier donné est contractile ;
  de façon équivalente, le groupe des automorphismes d'un collier qui
  induisent l'identité sur $B$ est contractile.
\end{enumerate}
Autrement dit, choisir un collier strict, ou un collier strict rigidifié, ou
une rigidification d'un collier donné, est anodin.

Pour des variétés topologiques à bord, a) est le théorème du collier de
M.~Brown (1962), et l'unicité des colliers à isotopie près est connue. Le
point c) demande une remarque, car lu à la lettre il ne peut pas valoir en
toute généralité. Deux rigidifications d'un même collier diffèrent d'un
homéomorphisme de $B \times I$ qui fixe $B \times \{0\}$ point par point et
préserve $B \times \{1\}$ ; pour $B$ une variété close, le groupe de ces
homéomorphismes est l'\emph{espace des pseudo-isotopies} (ou des
concordances) de $B$. Or cet espace n'est pas contractible en général : déjà
son $\pi_0$ est non nul pour certaines variétés compactes non simplement
connexes de grande dimension (Hatcher et Wagoner, 1973, dans le cadre lisse ;
l'obstruction vaut aussi dans le cadre linéaire par morceaux). Une variété
linéaire par morceaux compacte étant un espace modéré, l'énoncé c) tel quel
ne peut donc pas valoir pour de tels $B$.\footnote{Une incise de la page
dit les automorphismes « \emph{libres} » sur $\dot C$, lecture incertaine ;
et une version biffée parlait de colliers « rigidifiés stricts et fibrés
sur » \ldots. Il est possible qu'il ait eu en vue une notion plus restreinte
d'automorphisme — par exemple compatible à la projection sur $B$ — pour
laquelle la contractibilité est vraie et facile. La lecture ne choisit pas :
elle constate que la formulation écrite tombe sous l'obstruction des
pseudo-isotopies. Ce constat est de la lecture, non de la page, et mérite
l'examen d'un spécialiste.} Ce que l'obstruction laisse intact est l'unicité
des colliers à isotopie près, qui suffit à beaucoup d'usages.

\pagerange{148}{148}
\subsection*{II. Cônes autour d'un point}

Soit $a \in X$. Un \emph{cône} autour de $a$ est un couple
$\dot C \subsetneq C$, $a \in C \subset X$, avec un homéomorphisme modéré
\[
  \rho \colon (C, \dot C) \xrightarrow{\ \sim\ } \bigl(\mathrm{Cone}(Y), Y\bigr),
  \qquad \rho(a) = o,
\]
où $Y$ est un espace modéré, $\mathrm{Cone}(Y) = Y \times I / Y \times
\{0\}$ et $o$ le sommet. Il est \emph{strict} si $\dot C$ est bordant dans
$X \smallsetminus (C \smallsetminus \dot C)$, \emph{rigidifié} si l'on retient
$\rho$.

\textbf{Conjecture 2.}
\begin{enumerate}
\item[a)] Il existe un cône strict autour de tout point.
\item[b)] L'espace des cônes stricts autour de $a$ est contractile.
\item[c)] L'espace des rigidifications d'un cône donné est contractile,
  c'est-à-dire : le groupe des automorphismes de $\mathrm{Cone}(Y)$ qui
  induisent l'identité sur $Y$ et fixent $o$ est contractile.
\end{enumerate}

Ici la situation est meilleure qu'au paragraphe I. Le point a) est le
\emph{théorème de structure conique locale} : dans une structure o-minimale,
tout point d'un ensemble définissable a un voisinage définissablement
homéomorphe au cône sur son \emph{entrelacs} (Milnor, 1968, pour les
ensembles algébriques ; Bochnak--Coste--Roy pour les semi-algébriques ;
Coste et van den Dries dans le cadre o-minimal). Le point c) est vrai, par le
\emph{truc d'Alexander} : si $h$ est un tel automorphisme, on pose, pour
$t \in \,]0,1]$,
\[
  h_t(x) = t \cdot h(x / t) \ \text{ si } |x| \leqslant t, \qquad
  h_t(x) = x \ \text{ sinon},
\]
où $|x|$ désigne la coordonnée radiale et $x \mapsto t\cdot x$
l'homothétie du cône ; $h_1 = h$, $h_t$ tend vers l'identité quand $t \to 0$
parce que $h$ fixe $o$, et la formule recolle sur $|x| = t$ parce que $h$
fixe $Y$. Le chemin $t \mapsto h_t$ contracte le groupe sur l'identité, et
l'homothétie étant modérée (semi-algébrique sur un cône linéaire), la
contraction reste dans le cadre modéré.\footnote{La page porte
« (Alexandroff ?) ». Le procédé est celui de J.~W.~Alexander (1923), non de
P.~S.~Alexandroff ; le point d'interrogation est le sien.} Le contraste entre
I~c) et II~c) est instructif : le sommet du cône absorbe tout, le bord d'un
collier n'absorbe rien.

\pagerange{148}{149}
\subsection*{III. Le système des onze équivalences}

Il introduit six catégories d'espaces modérés, dont il ne précise pas les
flèches\footnote{Vraisemblablement des homéomorphismes modérés, avec pour
morphismes supérieurs les isotopies : c'est ce que dit l'adjectif
« $\infty$-isotopique ».} :
\begin{itemize}
\item $\mathrm{PB}$, les paires bordantes $(X, B)$ ; $\mathrm{PBC}$, les
  paires bordantes munies d'un collier strict $(C, \dot C)$ ;
  $\mathrm{PBCR}$, celles munies d'un collier strict rigidifié $(C,
  \varphi)$ ;
\item $\mathrm{EP}$, les espaces pointés $(X, a)$ ; $\mathrm{EPC}$, les
  espaces pointés munis d'un cône strict ; $\mathrm{EPCR}$, ceux munis d'un
  cône strict rigidifié.
\end{itemize}

Il affirme une « cascade » de foncteurs qui seraient tous des équivalences
de catégories $\infty$-isotopiques — en termes actuels, des équivalences
d'$\infty$-catégories, les espaces de morphismes étant des espaces
d'homéomorphismes :

\begin{tikzcd}[column sep=6em, row sep=large, nodes={font=\scriptsize}]
  & \mathrm{PB} \arrow[dl, "{\text{recoller } B \times I}"'] \arrow[dr, "{\text{recoller Cône}(B)}"] \arrow[dddl, bend right=20, "{\mathrm{id}}" description] \arrow[dddr, bend left=20] & \\
  \mathrm{PBCR} \arrow[rr, dashed] \arrow[d, "\text{oubli rig.}"'] & & \mathrm{EPCR} \arrow[d, "\text{oubli rig.}"] \\
  \mathrm{PBC} \arrow[rr, dashed, "{\text{contracter } B}"] \arrow[d, "\text{oubli collier}"'] & & \mathrm{EPC} \arrow[d, "\text{oubli cône}"] \\
  \mathrm{PB} \arrow[rr, "{(X,B) \mapsto (X/B, *)}"'] & & \mathrm{EP}
\end{tikzcd}

Les onze foncteurs sont les onze flèches du diagramme.\footnote{La page note
en marge « Système des 11 foncteurs qui sont des équivalences
$\infty$-isotopiques » ; le compte des flèches du diagramme transcrit donne
bien onze. L'étiquette de la flèche du bas, $(X, B) \mapsto (X/B, *)$, est
la nôtre : la page la laisse nue, et c'est la contraction de $B$ en un point,
comme la flèche horizontale médiane qui porte, elle, « contraction de $B$ en
un point ». Les deux 2-flèches $\alpha$ et $\beta$ que la page dessine entre
l'identité et les composés latéraux ne sont pas reproduites ; une marge dit
que pour les définir il faut « compliquer » les données, mots en partie
illisibles.} « Recoller $B \times I$ » remplace $(X, B)$ par
$(X \cup_{B} B \times I,\ B \times \{0\})$ : on ajoute un collier extérieur,
qui vient avec sa rigidification ; « recoller $\mathrm{Cone}(B)$ » remplace
$(X, B)$ par l'espace pointé $(X \cup_B \mathrm{Cone}(B), o)$, avec son cône
canonique. La flèche du bas contracte $B$ en un point ; son inverse attendu
est l'\emph{éclatement réel} du point : retirer l'intérieur d'un cône strict
et garder son entrelacs comme bord.

La lecture insiste sur un point que la page laisse implicite : ces foncteurs
ne peuvent être des équivalences qu'\emph{à isotopie près}, et seulement dans
un cadre modéré. Dans le cadre topologique, l'homéomorphisme
$(r, \theta) \mapsto (r, \theta + 1/r)$ du plan, qui fixe l'origine, ne se
relève à aucun homéomorphisme du plan éclaté en l'origine ; c'est une
spirale, et une structure o-minimale exclut précisément les spirales. C'est
une bonne illustration de ce que la topologie modérée est faite pour éviter.

Il note une variante : multipaires bordantes $B = \coprod_{i \in I} B_i$ et
espaces multipointés $I \hookrightarrow X$, $I$ fini variable. La page redessine
le système en perspective, et se termine sur une phrase en partie illisible
dont le sens est clair : ces catégories étant $\infty$-isotopiquement
équivalentes, on peut travailler avec n'importe laquelle.

\pagerange{150}{150}
\subsection*{IV. Voisinages tubulaires}

Soit $Y$ une partie modérée fermée d'un espace modéré $X$. Un
\emph{voisinage tubulaire} de $Y$ est par définition un cône strict autour du
point $o$ image de $Y$ dans $X/Y$ — c'est-à-dire, en remontant, un
voisinage $V$ de $Y$ dont le « bord » $\dot V$ est bordant et tel que
$V/Y$ soit un cône sur $\dot V$.

\emph{Construction.} Soit $f \colon X \to \mathbb{R}_{\geqslant 0}$ une
fonction modérée telle que $f^{-1}(0) = Y$ ; pour $\varepsilon > 0$ assez
petit, on prend
\[
  V = f^{-1}([0, \varepsilon]), \qquad \dot V = f^{-1}(\varepsilon).
\]
C'est la construction qu'on appelle aujourd'hui voisinage défini par une
\emph{fonction tapis} (« rug function ») ; pour les ensembles semi-algébriques,
Durfee (1983) a montré que le résultat, pour $\varepsilon$ petit, ne dépend
pas de $f$ à isotopie près.\footnote{Le « petit » est accompagné sur la
page d'un « [et admissible] » qu'il ne précise pas. Rien ne dit qu'il
connaissait l'article de Durfee, \emph{Neighborhoods of algebraic sets},
Trans. AMS 276 (1983).}

\emph{Rigidification.} Rigidifier le voisinage tubulaire, c'est rigidifier le
cône, ce qui revient à se donner
\[
  V \smallsetminus Y \;\simeq\; \dot V \times [0, 1[
\]
et une application modérée $\varphi \colon \dot V \to Y$ telle que
\[
  X \;\simeq\; \Bigl( \bigl(X \smallsetminus (V \smallsetminus \dot V)\bigr)
  \cup_{\dot V \times \{0\}} (\dot V \times I) \Bigr)
  \cup_{\dot V \times \{1\},\, \varphi} Y .
\]
Autrement dit, $V$ est le \emph{cylindre de l'application} $\varphi$, recollé
au reste de $X$ le long de $\dot V$ : c'est un voisinage en cylindre
d'application.\footnote{La page écrit $X \smallsetminus (\dot V \smallsetminus
V)$, ce qui est vide de sens ($\dot V \subset V$) ; la lecture rétablit
$X \smallsetminus (V \smallsetminus \dot V)$, l'extérieur de $V$ bord compris,
seul compatible avec le recollement. Il dit $\varphi$ « déterminée de façon
unique par la rigidification ».} Il appelle cette donnée une rigidification
\emph{$Y$-admissible}, et énonce :
\begin{enumerate}
\item[a)] une rigidification $Y$-admissible existe ;
\item[b)] la classe d'homotopie de $\varphi \colon \dot V \to Y$ est bien
  déterminée ;
\end{enumerate}
« mais on voudrait plus précis », et en marge : « choix de voisinages
tubulaires et de rigidifications anodins ».

Ces deux points sont aujourd'hui acquis dans le cadre o-minimal compact. Pour
a), on triangule $X$ de façon définissable avec $Y$ sous-complexe ; le voisinage simplicial de $Y$ dans la deuxième subdivision barycentrique est,
par sa structure de joint, homéomorphe au cylindre d'une application de sa
frontière vers $Y$. Pour b),
$V$ se rétracte par déformation sur $Y$, et $\varphi$ est homotope à
l'inclusion $\dot V \hookrightarrow V$ suivie de cette rétraction : sa classe
d'homotopie ne dépend que de l'inclusion. C'est le « plus précis » — un
énoncé de contractibilité des choix — qui reste au niveau d'un programme.

\pagerange{151}{152}
\subsection*{V. Le cas équisingulier, le cas général, les filtrations}

\emph{1° Cas où $X$ est équisingulier le long de $Y$.} Il présume qu'on peut
alors prendre $\varphi$ \emph{fibrante} — une fibration localement triviale
— et que l'on a une équivalence de catégories $\infty$-isotopiques
\[
  \bigl\{ (X', B, Y, \varphi) \bigr\}
  \;\xrightarrow{\ \approx\ }\;
  \bigl\{ (X, Y) \text{ paires équisingulières} \bigr\},
  \qquad
  (X', B, Y, \varphi) \longmapsto \bigl( X' \cup_{B, \varphi} Y,\ Y \bigr),
\]
où $(X', B)$ est une paire bordante et $\varphi \colon B \to Y$ une
fibration. Quand $X \smallsetminus Y$ et $Y$ sont réguliers (des variétés),
cela doit correspondre au cas où $X'$ est une variété à bord $B$.

C'est, sous un autre nom, la description d'un espace stratifié à deux strates
par sa \emph{résolution} : une variété à bord dont le bord est fibré sur la
strate fermée, la fibre étant l'entrelacs. Dans le cadre des espaces
stratifiés de Whitney, l'existence de la fibration est le premier lemme
d'isotopie de Thom et la théorie des données de contrôle de Mather (1970) ;
la version « équivalence de catégories » de son énoncé est très proche du
« dézippage » d'Ayala, Francis et Tanaka (2017) pour les espaces stratifiés
coniquement lisses.\footnote{Rapprochement de la lecture : la page ne nomme
ni Thom ni Mather, et le travail d'Ayala--Francis--Tanaka lui est très
postérieur. La marge de la page 151 contient une condition sur $Y$ (« $Y
\subset \overline{X - Y}$ ») entourée de mots illisibles.}

\emph{2° Cas général.} Sans équisingularité, il faudrait imposer à
$\varphi \colon \dot V \to Y$ des conditions assurant l'unicité essentielle,
par exemple une condition de dimension : si $c(y)$ est la codimension de $Y$
dans $X$ en $y$, alors
\[
  \dim \varphi^{-1}(y) = c(y) - 1,
\]
comme pour le fibré en sphères normal dans le cas lisse. Il « rêve » d'une
description de la catégorie $\infty$-isotopique des paires $(X, Y)$ en termes
de données $(X', B, \varphi \colon B \to Y)$ soumises à de telles
conditions.\footnote{Le mot qualifiant les paires est lu « emplies », très
incertain ; la phrase sur les conditions locales de $\dot V$ au-dessus de
$Y$ se termine sur un mot illisible.} Les « voisinages en goutte »
(\emph{teardrop neighbourhoods}) de Hughes et les espaces homotopiquement
stratifiés de Quinn (1988) sont la forme la plus générale connue de ce rêve
dans le cadre topologique.

\emph{Filtrations.} Après un trait, une dernière question : un espace modéré
a-t-il une filtration \emph{canonique}, régulière et équidimensionnelle ? Il
le présume, et note que les invariants de dévissage de cette filtration
seraient alors des invariants de la structure modérée elle-même.\footnote{Le
passage entre crochets qui précise l'énoncé est en partie biffé et
illisible.} La réponse est oui sous des formes diverses : les ensembles
semi-algébriques ont une stratification de Whitney minimale, et les
pseudo-variétés une stratification intrinsèque, purement topologique (King,
1985), dont l'homologie d'intersection est un invariant.

\pagerange{154}{160}
\section*{Digressions combinatoires (pages 154 à 160)}

Sept pages photocopiées d'un manuscrit, sous le titre « Digressions
combinatoires (pour Réflexions t.~4) ».\footnote{On n'identifie pas ici avec
certitude le « tome 4 » de « Réflexions » auquel il destinait ce texte.}
Ce n'est pas une rédaction mais un plan, plusieurs fois repris. Le sujet est
une géométrie des petits ensembles finis, qui rend visibles les
\emph{isomorphismes exceptionnels} entre groupes symétriques et groupes
linéaires ou affines sur de petits corps finis, et l'\emph{automorphisme
extérieur} de $\mathfrak{S}_6$.

\subsection*{Deux notions de base}

\emph{Orientations et produit contracté.} Un ensemble à deux éléments sur
lequel agit $\mathbb{Z}/2$ est un $\mathbb{Z}/2$-torseur. Si $\varepsilon$ et
$\varepsilon'$ en sont deux, leur \emph{produit contracté}
$\varepsilon \wedge \varepsilon'$ est le quotient de $\varepsilon \times
\varepsilon'$ par l'action diagonale ; concrètement, c'est l'ensemble des deux
bijections $\varepsilon \to \varepsilon'$. On a canoniquement
$\varepsilon \wedge \varepsilon \simeq 1$ (le torseur trivial). Une
\emph{orientation} d'un ensemble fini $S$ (de cardinal $\geqslant 2$) est une
numérotation de ses éléments à permutation paire près ; l'ensemble $\omega(S)$
de ses deux orientations est un torseur, et l'on a
$\omega(S \amalg T) \simeq \omega(S) \wedge \omega(T)$, ainsi que
$\omega(\varepsilon) \simeq \varepsilon$ pour un ensemble $\varepsilon$ à deux
éléments. Ses calculs des pages 155 et 156 sont des identités entre de tels
torseurs.

\emph{Carrés.} Une structure de \emph{carré} sur un ensemble $E$ à quatre
éléments est une partition de $E$ en deux paires — les deux diagonales du
carré. Il y en a trois. Il note $\mathrm{Car}(E)$ leur ensemble et l'identifie
à l'ensemble des « antipodismes » de $E$.

\pagerange{154}{154}
\subsection*{Le premier plan (page 154)}

Le plan en quatre chapitres est : géométrie des ensembles de cardinal
$\leqslant 4$ ; l'icosaèdre « gauche » et ses trois perspectives ; les
bi-icosaèdres ; la dualité des hexagrammes. Ce dernier chapitre est le cœur,
et voici ce qu'il dit en termes actuels.

Soit $S$ un ensemble à six éléments. Les sous-groupes de $\mathfrak{S}(S)$
isomorphes à $\mathfrak{S}_5$ et opérant transitivement sur $S$ sont au nombre
de six ; leur ensemble est un nouvel ensemble à six éléments $S^*$, le
\emph{dual} de $S$, et l'on a $S^{**} \simeq S$. L'isomorphisme
$\mathrm{Aut}(S) \simeq \mathrm{Aut}(S^*)$ qui en résulte est l'automorphisme
extérieur de $\mathfrak{S}_6$ : c'est ce qu'il appelle l'\emph{autodualité}
$\mathrm{Ens}_6 \simeq \mathrm{Ens}_6$.\footnote{Description classique,
par les « synthèmes » et « totaux » de Sylvester (1844). La page ne nomme
personne.}

Il nomme \emph{bi-icosaèdre} ce qui est fixé par un tel $\mathfrak{S}_5$
transitif. L'interprétation est la suivante : un icosaèdre a six axes
(droites joignant deux sommets opposés), et son groupe des rotations,
isomorphe à $\mathfrak{A}_5$, opère sur ces six axes comme
$\mathrm{PSL}_2(\mathbb{F}_5)$ sur les six points de la droite projective
$\mathbf{P}^1(\mathbb{F}_5)$ ; le normalisateur
$\mathrm{PGL}_2(\mathbb{F}_5) \simeq \mathfrak{S}_5$ échange les deux
icosaèdres « miroirs » qui ont les mêmes axes. Une structure
\emph{icosaédrale} sur $S$ (identification à l'ensemble des axes d'un
icosaèdre) a donc pour groupe $\mathfrak{A}_5$ ; il y en a douze, groupées en
six paires — les six bi-icosaèdres, c'est-à-dire les six points de $S^*$. Les
éléments qui échangent les deux icosaèdres d'une paire agissent sur $S$ par
des permutations impaires\footnote{Par exemple $x \mapsto 2x$ sur
$\mathbf{P}^1(\mathbb{F}_5)$ fixe $0$ et $\infty$ et permute $1, 2, 4, 3$
circulairement : c'est un 4-cycle, impair.} ; c'est ce qui donne sens à sa
formule
\[
  \mathrm{Ic}(S) \;\simeq\; S^* \times \omega(S) ,
\]
bijection compatible à l'action de $\mathfrak{S}(S)$.\footnote{La page
n'indique pas laquelle des deux bijections équivariantes possibles est prise :
il y en a exactement deux, échangées par l'involution de $\omega(S)$. Le
choix est une convention globale.}

Le « dictionnaire de la dualité » apparie les structures sur $S$ et sur
$S^*$ selon les types de partitions de six, ce qui est la façon dont
l'automorphisme extérieur agit sur les classes de conjugaison de sous-groupes :
\begin{itemize}
\item un point marqué de $S$ (type $(1,5)$) correspond à un bi-icosaèdre sur
  $S^*$, et un point marqué avec une orientation à un icosaèdre ;
\item une paire marquée (type $(2,4)$, groupe $\mathfrak{S}_2 \times
  \mathfrak{S}_4$) correspond à une partition en trois paires (type $(2,2,2)$,
  groupe $\mathfrak{S}_2 \wr \mathfrak{S}_3$), c'est-à-dire à une structure de
  cube, les trois paires étant les trois paires de faces opposées ; les deux
  groupes sont d'ordre $48$ ;
\item une partition en deux triades (type $(3,3)$, « ditriade », groupe
  $\mathfrak{S}_3 \wr \mathfrak{S}_2$ d'ordre $72$) correspond à une
  ditriade.
\end{itemize}
Une seconde colonne (A$'$, B$'$, C$'$) ajoute à chaque structure un sommet,
une arête ou une « biface », et la page 156 en tire les correspondances
fines.\footnote{La page 154 comporte un premier jet de cette seconde colonne
entièrement biffé, et plusieurs ratures dans le reste ; la lecture ne reprend
que ce qui est lisible et non biffé.}

\pagerange{155}{156}
\subsection*{Carrés duaux et plan affine sur $\mathbb{F}_3$ (pages 155 et 156)}

\emph{Les carrés.} Une structure de carré $Q$ sur un sous-ensemble à quatre
éléments de $S$ correspond, par la dualité, à une structure de carré $Q^*$
sur un sous-ensemble à quatre éléments de $S^*$, et il écrit les relations
\[
  \varepsilon \wedge \varepsilon^* \wedge \omega(Q) \simeq 1, \qquad
  \omega(Q) \simeq \omega(Q^*), \qquad
  Q^* \simeq \varepsilon \wedge Q, \qquad Q \simeq \varepsilon^* \wedge Q^*,
\]
où $\varepsilon$, $\varepsilon^*$ sont les paires complémentaires.\footnote{La
page 156 encadre la même relation sous la forme $Q^* = \varepsilon \wedge
Q^{\vee}$, avec un $Q^{\vee}$ qui n'est pas défini ailleurs ; la page 155
écrit $Q^* \simeq \varepsilon \wedge Q$. La lecture suit la page 155. Le sens
précis de « $\varepsilon \wedge Q$ » (le produit contracté d'un torseur avec un
ensemble de carrés) n'est pas défini sur les pages.} La page 156 les combine
avec les torseurs $\mathrm{diag}\,Q$ et $\mathrm{codiag}\,Q$ et aboutit à
$\mathrm{codiag}\,Q^* \simeq \mathrm{diag}\,Q$, puis, dans un « programme
icosaèdre », lit ces torseurs sur les sommets d'un icosaèdre dessiné en
perspective. Elle note enfin, avec un « OK ! », la vérification
\[
  \omega(S \amalg \varepsilon) \simeq \omega(S) \wedge \varepsilon
  \simeq \mathrm{codiag}(Q) \wedge \varepsilon \simeq \mathrm{diag}\,Q,
\]
qui identifie le choix d'une structure icosaédrale dans un bi-icosaèdre à
$\mathrm{diag}\,Q$. Ces calculs dépendent de définitions que les pages ne
donnent pas ; la lecture les rapporte sans les refaire.\footnote{Page de
calculs serrée, disposée en tableau, avec des renvois entourés (1), (2),
(3), (1$'$)\ldots{} aux correspondances de la page 154 et plusieurs
ratures.}

La même page 156 note deux faits justes sur $\mathfrak{S}_6$ : ses 3-sous-groupes
de Sylow sont isomorphes à $\mathbb{F}_3^2$ et ses 2-sous-groupes de Sylow à
$D_4 \times \mathbb{Z}/2$ (ordres $9$ et $16$).

\emph{Le plan affine sur $\mathbb{F}_3$.} Soit $V$ un plan affine sur
$\mathbb{F}_3$ (neuf points). L'ensemble $E = \mathbf{P}(\vec V)$ de ses
directions a quatre éléments, et
\[
  \mathfrak{S}(E) \simeq \mathrm{PGL}_2(\mathbb{F}_3),
\]
le noyau de $\mathrm{GL}_2(\mathbb{F}_3) \to \mathrm{PGL}_2(\mathbb{F}_3)$
étant $\{\pm 1\}$ — « l'antipodisme ». Pour chaque direction $i \in E$,
l'ensemble $E_i$ des droites de direction $i$ a trois éléments, $V$ se plonge
dans $\prod_i E_i$, et deux directions distinctes donnent des coordonnées :
$V \xrightarrow{\sim} E_i \times E_j$ pour $i \neq j$. Le groupe des
automorphismes de cette structure — il l'appelle « trilitaire » — est
$\mathrm{AGL}_2(\mathbb{F}_3)$, d'ordre $3^3 \cdot 2^4 = 432 = 3 \cdot 144$,
comme il l'écrit. Il opère sur l'ensemble $\coprod_{i} E_i$ des douze droites
et permute trois sous-groupes $\mathfrak{S}_\alpha$ indexés par les trois
carrés $\alpha \in \mathrm{Car}(E)$ ; le stabilisateur d'un $\mathfrak{S}_\alpha$
est le groupe des automorphismes d'une \emph{bi-ditriade}.\footnote{La page
ajoute que ce stabilisateur opère fidèlement sur $\mathfrak{S}_\alpha$ comme
« sous-groupe d'ordre 10 » de $\mathrm{Aut}(\mathfrak{S}_\alpha)$, d'ordre
$2 \cdot (3!)^2 = 72$ (lecture incertaine). Comme lu, c'est impossible :
$10$ ne divise pas $72$. La lecture ne reconstruit pas le nombre ; elle
signale aussi que la définition des $\mathfrak{S}_\alpha$ n'est pas sur la
page, et que la note marginale qui semble en parler (« 3-sous-groupes
\ldots{} Sylow ») est presque illisible.}

\pagerange{157}{157}
\subsection*{Une feuille intercalée (page 157)}

La page 157 n'appartient pas au plan. C'est une liste dactylographiée de
couples de mots, renumérotée à sa main, et qui se lit comme une table de
couples complémentaires : « Avenir — Passé », « destinée — histoire »,
« pérennité — ancienneté », « innovation — tradition », « élan —
enracinement » sous la rubrique ajoutée « Devenir » ; puis, sous
« Espace-temps », « Espace — temps » et « étendue (ou distance) —
durée ».\footnote{Plusieurs couples sont biffés (« le nouveau — l'ancien »,
« ubiquité — éternité »), et les numéros de rubrique sont corrigés de
IV$''_5$ en V$''_4$ et de IV$''_6$ en V$''_5$.} Viennent ensuite des projets
de notes de bas de page numérotées 98 à 107 pour un texte de lui, qui
renvoient à une « note n°~162 » et à une « p.~779 » (« conviction et
connaissance »), à « plusieurs existences successives », à un « cours 1977 »,
et des croquis de tétraèdres et d'octaèdres avec une liste de faces.
\footnote{Le texte visé n'est pas identifié ici ; plusieurs mots sont
illisibles.} Elle est de sa main et de sa pensée, et se trouve là par le
hasard de la photocopie.

\pagerange{158}{158}
\subsection*{Trialité des bi-ditriades (page 158)}

La page décrit une même structure de deux façons, par les neuf points du plan
affine $V$ sur $\mathbb{F}_3$ ou par ses douze droites $E$, liés par
l'incidence.

\emph{A) Sur $V$.} La structure revient à la donnée d'un ensemble $E$ de
douze parties à trois éléments de $V$, deux à deux d'intersection au plus un
point. C'est exact, et la raison est un comptage : douze triplets contiennent
$12 \times 3 = 36$ paires de points, qui est le nombre $\binom{9}{2}$ de toutes
les paires ; comme deux triplets n'ont jamais deux points communs, chaque paire
est dans exactement un triplet. On a un \emph{système triple de Steiner} à
neuf points, et il n'y en a qu'un à isomorphisme près : le plan affine
$\mathrm{AG}(2, 3)$.\footnote{La justification par comptage est de la
lecture ; la page énonce sans démontrer.}

\emph{B) Sur $E$.} La structure revient à la donnée d'une partition
$E = \coprod_{i \in I} E_i$ en quatre classes de trois éléments (les quatre
directions), et d'un ensemble $V$ de parties à quatre éléments de $E$ (les
quatre droites passant par un point) telles que :
\begin{enumerate}
  \item[a)] chaque $D \in V$ rencontre chaque $E_i$ en exactement un élément ;
  \item[b)] si $D \neq D'$ dans $V$, alors $\mathrm{card}(D \cap D') = 1$ ;
  \item[c)] si $a, b \in E$ ne sont pas dans la même classe, il existe un
    unique $D \in V$ qui les contient.
\end{enumerate}
Il note que b) résulte de a) et c), et que c) implique $\mathrm{card}\,V = 9$,
puis, compte tenu de a), que par tout $a \in E$ passent exactement trois
éléments de $V$. Le comptage est le suivant : il y a
$\binom{12}{2} - 4\binom{3}{2} = 54$ paires d'éléments de classes distinctes,
chaque $D$ en contient $\binom{4}{2} = 6$, et chaque paire est dans un seul
$D$ ; d'où $\mathrm{card}\,V = 9$.\footnote{Ce comptage utilise que chaque
$D$ a quatre éléments pris dans des classes distinctes, c'est-à-dire a) :
c) seul n'implique pas $\mathrm{card}\,V = 9$, contrairement à ce que dit la
page.} De même, $a$ forme une paire avec chacun des neuf éléments des autres
classes, et chaque $D$ contenant $a$ en fournit trois : trois droites
par $a$.

Une marge dit l'énoncé « dans un plan projectif sur $\mathbb{F}_3$ privé d'un
point ». C'est exact par dualité : $\mathrm{AG}(2,3)$ est le plan projectif
$\mathbf{P}^2(\mathbb{F}_3)$ privé d'une droite, et le dual d'une droite est
un point ; les douze éléments de $E$ sont les douze points de
$\mathbf{P}^2(\mathbb{F}_3)$ autres que ce point, les quatre classes sont les
quatre droites qui y passent, et les neuf éléments de $V$ sont les neuf droites
qui n'y passent pas.\footnote{Une autre marge, en biais, en grande partie
illisible, commence par « $E \cong V \amalg \{\infty\}$ ».}

\pagerange{159}{160}
\subsection*{Le second plan : géométrie des ensembles finis (pages 159 et 160)}

Le plan est repris en cinq parties, une par cardinal. Les isomorphismes qu'il
énonce sont tous exacts ; la lecture les donne en notation actuelle.

\emph{I. Deux éléments.} $\mathfrak{S}_2 \simeq \{\pm 1\} \simeq
\mathbb{Z}/2 \simeq \mathbb{F}_3^{*}$ ; le produit contracté ; les orientations
d'un ensemble fini et leur lien avec l'orientation d'un espace vectoriel réel
(signe du déterminant), d'un espace orthogonal en caractéristique $\neq 2$, ou
d'un espace vectoriel sur $\mathbb{F}_3$ (où le déterminant d'une base à une
autre est dans $\mathbb{F}_3^{*} = \{\pm1\}$) ; hypercubes et antipodismes ;
carrés.

\emph{II. Trois éléments.} $\mathfrak{S}_3 \simeq \mathrm{GL}_2(\mathbb{F}_2)
= \mathrm{SL}_2(\mathbb{F}_2) = \mathrm{PGL}_2(\mathbb{F}_2)$, par l'action sur
les trois points de $\mathbf{P}^1(\mathbb{F}_2)$ ; et $\mathfrak{S}_3 \simeq
\mathrm{AGL}_1(\mathbb{F}_3)$, par l'action sur les trois points de la droite
affine, le passage à la partie linéaire $\mathrm{AGL}_1(\mathbb{F}_3) \to
\mathbb{F}_3^{*} = \{\pm 1\}$ étant la signature. Application aux ditriades et
bi-ditriades, avec une suite exacte
$1 \to \mathfrak{S}_3 \times \mathfrak{S}_3 \to \mathrm{Aut}(\text{bi-ditriade})
\to \cdots \to 1$ dont le dernier terme est biffé et illisible.

\emph{III. Quatre éléments.} $\mathfrak{S}_4 \simeq
\mathrm{AGL}_2(\mathbb{F}_2)$ (les quatre points du plan affine sur
$\mathbb{F}_2$) ; $\mathfrak{S}_4 \simeq$ groupe des rotations du cube, et une
structure de cube revient à la donnée d'un ensemble à quatre éléments et d'un
ensemble à deux ($\mathfrak{S}_4 \times \mathfrak{S}_2$, d'ordre 48) ;
$\mathfrak{S}_4 \simeq \mathrm{PGL}_2(\mathbb{F}_3)$ (les quatre points de
$\mathbf{P}^1(\mathbb{F}_3)$).\footnote{La fin de la ligne sur le cube est une
addition marginale serrée de lecture très incertaine ; la lecture ne s'en
sert pas.} Enfin les \emph{quadritriades}, dont le groupe d'automorphismes
est
\[
  \mathrm{Aut}(\text{quadritriade}) \simeq \mathrm{AGL}_2(\mathbb{F}_3),
  \qquad
  1 \to (\mathfrak{S}_3 \times \mathfrak{S}_3)' \to \mathrm{AGL}_2(\mathbb{F}_3)
  \to \mathfrak{S}_4 \to 1,
\]
où $(\mathfrak{S}_3 \times \mathfrak{S}_3)' = \{ (u, v) \mid
\mathrm{sgn}(u) = \mathrm{sgn}(v) \}$. On le vérifie : la flèche de droite est
l'action sur les quatre directions, $\mathrm{AGL}_2(\mathbb{F}_3) \to
\mathrm{PGL}_2(\mathbb{F}_3) \simeq \mathfrak{S}_4$ ; son noyau est formé des
translations et de $\pm 1$, groupe d'ordre $18$ isomorphe à
$\mathbb{F}_3^2 \rtimes \{\pm 1\}$, qui est bien $(\mathfrak{S}_3 \times
\mathfrak{S}_3)'$ ; et $18 \times 24 = 432$.

\emph{IV. Cinq éléments.} $\mathfrak{A}_5 \simeq
\mathrm{PGL}_2(\mathbb{F}_4)$, par l'action sur les cinq points de
$\mathbf{P}^1(\mathbb{F}_4)$ ; et $\mathfrak{S}_5$ est le groupe des
automorphismes d'une droite projective sur un corps à quatre éléments « non
précisé » — c'est-à-dire $\mathrm{P\Gamma L}_2(\mathbb{F}_4)$, le Frobenius
compris, ce que dit exactement « non précisé ». Il ajoute que si $E =
\mathbf{P}^1(k)$ avec $k \simeq \mathbb{F}_4$, alors $\omega(E) \simeq k^{*}
\smallsetminus \{1\}$, l'ensemble des deux générateurs de $k$ sur
$\mathbb{F}_2$ : c'est juste, le Frobenius échangeant ces deux générateurs et
agissant sur les cinq points par une transposition. Enfin $\mathfrak{S}_5
\simeq \mathrm{PGL}_2(\mathbb{F}_5)$, par l'action sur les six points de
$\mathbf{P}^1(\mathbb{F}_5)$ — c'est le $\mathfrak{S}_5$ transitif de la page
154, et il le renvoie au bi-icosaèdre.

\emph{V. Six éléments} et bi-icosaèdres, « voir plus détaillé ailleurs » ; puis
un appendice sur les polygones combinatoires, et le constat qu'il faudrait un
exposé sur les relations géométriques entre toutes ces configurations. Le
plan s'arrête là.

\pagerange{241}{250}
\section*{« Le maître-enseignant et le maître-chercheur » (juin 1968)}

\subsection*{Le contexte}

Le texte est écrit au début de juin 1968, au milieu de la grève générale et
du mouvement étudiant de mai. Grothendieck est alors professeur à l'Institut
des Hautes Études Scientifiques, à Bures-sur-Yvette. Le département de
mathématiques de la Faculté des sciences d'Orsay — qui deviendra en 1971
l'université Paris-Sud — a constitué une commission sur la formation et le
recrutement des enseignants et chercheurs ; ses collègues l'y ont invité,
comme le dit sa lettre au \emph{Monde}.

Le vocabulaire des grades a changé depuis, et le texte ne se comprend qu'avec
celui de 1968. L'\emph{assistant} et le \emph{maître-assistant} formaient le
corps enseignant subalterne, nombreux et non titulaire d'un doctorat
d'État ; le \emph{maître de conférences} et le \emph{professeur} formaient le
corps des « maîtres », peu nombreux, recruté sur la thèse de doctorat d'État,
c'est-à-dire sur la recherche.\footnote{Le titre a changé de sens depuis : en
1979 les maîtres de conférences d'alors sont devenus professeurs, et en 1984
le corps des maîtres-assistants a pris le nom de maîtres de conférences. Le
« maître de conférences » du texte est donc l'équivalent d'un professeur
d'aujourd'hui. Le doctorat d'État a été remplacé en 1984 par le doctorat
unique et l'habilitation à diriger des recherches.} La « licence » est le
diplôme de fin du premier cycle long ; l'enseignement « pré-licence » est
ce qu'on appelle aujourd'hui les premières années de licence.

Le dossier en conserve deux états, que la lecture présente ensemble. Le
\emph{premier état} (pages 241 à 250) est un tapuscrit de dix pages, titré
« \ldots{} dans l'Université d'aujourd'hui », auquel il a ajouté à la main
« et de demain » et l'intertitre « A) Aujourd'hui », et dont il a récrit la
fin pour annoncer une suite. C'est ce texte, sous son titre premier, qu'il
envoie au \emph{Monde} le 4 juin. Le \emph{second état} (pages 254 à 270)
est le même texte retapé, les ajouts incorporés, et augmenté d'une partie
« B. Demain » de huit points.\footnote{Les différences entre les deux états
sur la partie A sont minimes et relevées dans l'en-tête de la transcription :
« dans une large mesure » devient « dans une certaine mesure » (à propos de
l'insuffisance des crédits, pages 244 et 257) ; une note est ajoutée sur le
« bon enseignant » (page 259) ; et quelques fautes de frappe nouvelles
apparaissent. Les corrections à l'encre du premier état sont reprises.}

\subsection*{A. Aujourd'hui : un vice de structure}

\emph{1. La confusion des deux rôles.} Depuis ses origines, l'Université
confie à la même personne deux rôles « étroitement apparentés, mais
néanmoins essentiellement distincts » : transmettre le savoir acquis, et le
créer. N'était censé pouvoir enseigner la science que celui qui contribuait à
l'édifier. Ce système convenait tant que les étudiants étaient peu nombreux et
allaient tous jusqu'aux frontières de la recherche : il y avait alors plus de
savants que de chaires.

\emph{2. Ce qui a changé.} L'essor de la technique et l'afflux d'étudiants
ont changé à la fois le nombre et la nature des besoins. La plupart des
étudiants en mathématiques seront ingénieurs, physiciens, chimistes,
biologistes, économistes : des « utilisateurs » de mathématiques, pas des
mathématiciens. Il leur faut une plus grande variété d'enseignements —
variété « non seulement par le nombre des théories enseignées, mais par
l'optique dans laquelle elles sont enseignées », libérée d'exigences de
rigueur et d'exhaustivité « justifiées il y a cent ans ». Un grand chercheur,
absorbé par ses problèmes, enseigne mal à ces étudiants-là, et c'est un double
gaspillage.

Surtout, il n'y a pas assez de chercheurs créateurs pour fournir les
enseignants nécessaires. En maintenant le critère de recherche, on recrute
trop peu de maîtres (a) ; on les fait assister par une masse croissante
d'assistants et de maîtres-assistants subalternes, sans perspective de
promotion (b) ; et les amphithéâtres restent surpeuplés (c). Le point c),
qu'il impute largement au manque de crédits, est pour les étudiants le vice
principal et l'une des causes de « la salutaire crise de l'Université que nous
sommes en train de traverser ». Les points a) et b), eux, viennent de la
confusion des deux rôles, et du préjugé qui en découle contre l'enseignant qui
n'est « que » enseignant.

Ce préjugé fausse aussi les carrières. Des scientifiques cultivés, au jugement
sûr, excellents pédagogues, mais qui ne produisent pas d'œuvre originale, sont
écartés de l'enseignement supérieur. D'autres, moins scrupuleux, obtiennent
une thèse complaisante, deviennent maîtres de conférences, puis jugent du
recrutement des chercheurs et dirigent des thèses « soi-disant de recherche »,
au risque d'étouffer la vie scientifique de leur département.

\emph{3. Le maître-chercheur.} Il s'intéresse d'abord aux problèmes que lui
pose la science, « plus qu'aux hommes qui l'apprennent ou l'utilisent, plus
même qu'aux hommes qui la font ». L'enseignement lui assure une position
sociale. Il est souvent bon enseignant, d'autant meilleur que le niveau est
élevé, et il peut l'être au niveau le plus élémentaire devant de futurs
mathématiciens ; mais il arrive aussi qu'un grand chercheur soit un
enseignant exécrable, dont les cours restent « touffus et laborieux quand ils
ne sont obscurs et sibyllins ». Le second état ajoute en note que la
réputation de « bon enseignant » est celle qu'on a parmi ses pairs, et que
des étudiants trouvent parfois l'assistant meilleur que son patron.

Il en donne un signalement qui se lit aujourd'hui comme un autoportrait : le
chercheur mathématicien est celui qui, rencontrant un congénère « autour
d'une table, dans un métro, dans un congrès », se met aussitôt à lui parler
non de politique ni de ses collègues, mais de problèmes mathématiques ; et
l'enseignement le passionne par les problèmes d'exposition qu'il pose, non par
ses auditeurs. Il rapporte, sans en garantir l'authenticité, l'anecdote de
C.~L.~Siegel faisant cours devant une salle vide le jour où son unique
auditeur était absent.

\emph{4. Le maître-enseignant.} Il s'intéresse à la science « autant ou plus
par les liens humains qu'elle implique, crée ou développe » que par la
découverte. Il a souvent fait une thèse de recherche au prix d'un effort
unique dans sa vie, puis ne produit plus guère. Le danger qui le guette est de
devenir le \emph{maître-fossile}, qui répète jusqu'à la retraite le cours mis
au point dans sa jeunesse — « véritable honte de notre métier », mais
sous-produit des structures plus que coupable. Le vrai enseignant sait qu'il
doit se renouveler, et cela n'est possible qu'au contact de la science
vivante : non qu'il doive tout suivre, mais il doit rester intellectuellement
actif, et peut contribuer à la recherche dans la mesure de ses moyens, comme
moyen et non comme fin. Car « le sujet principal d'intérêt du
maître-enseignant est l'étudiant lui-même », ce qui suppose un contact
étroit, à l'opposé des « amphithéâtres-forums » où le
« maître-gardien-de-la-science » fait face à la foule comme dans une
« hallucinante vision de Kafka ».

\emph{5. Conclusion de la première partie.} Il ne s'agit pas d'établir une
hiérarchie entre les deux vocations, mais de réhabiliter la vocation
pédagogique en la mettant « à pied d'égalité » avec la recherche. Les étudiants
ont été les principales victimes du système ; il faut profiter de la « force
vive » libérée par leur mouvement pour mettre au point, dès la rentrée, des
mécanismes souples qui distinguent et coordonnent les deux vocations. La fin
est la même dans les deux états, une fois ses ratures intégrées : son but a été
de mettre en évidence « un des vices de base du système actuel », et
d'inviter chacun à y réfléchir, car c'est « de nous tous qui sommes
concernés, étudiants, enseignants et chercheurs, non des princes qui nous
gouvernent », que dépend la rénovation.\footnote{Dans le premier état, la fin
de la page 249 est récrite à la main : « Il n'est pas dans mon propos ici de
passer en revue » devient « Je passerai en revue plus bas », et la phrase qui
disait vouloir laisser les solutions « à telles assemblées spontanées plus
autorisées que moi » est biffée. C'est la trace du moment où le texte, de
diagnostic, devient programme : la partie B est annoncée à l'encre avant
d'être écrite. La transcription signale un « d' » surfrappé, peut-être sur
un « à ».}

\pagerange{251}{252}
\subsection*{La lettre à Hubert Beuve-Méry (4 juin 1968)}

Lettre dactylographiée, signée à la main, de Massy, adressée au directeur du
\emph{Monde}.\footnote{Hubert Beuve-Méry (1902--1989) a fondé \emph{Le
Monde} en 1944 et l'a dirigé jusqu'en 1969.} Grothendieck y explique qu'il
participe à la commission d'Orsay, que les discussions entre étudiants,
assistants et professeurs ont montré que certains vices fondamentaux du
système, clairs pour les professeurs, l'étaient moins pour les étudiants et les
assistants, pourtant directement concernés ; et qu'il faut informer un public
plus large — au moins tous les étudiants et enseignants de mathématiques, et
même de sciences, les problèmes étant sans doute les mêmes, quoique moins
aigus, dans les autres disciplines. Il joint son exposé, sous son titre
premier, en suggère la publication en un ou deux articles de fond, et demande
une réponse rapide pour pouvoir, sinon, le diffuser autrement. « A cause de la
grève des postes », il demande qu'on lui réponde par téléphone. Il signe
« Professeur à l'Institut des Hautes Études Scientifiques ».\footnote{Le
dossier ne dit pas si le journal a répondu ; le relevé du dossier (issue
\#45) n'a trouvé le texte imprimé nulle part.}

\pagerange{254}{262}
\subsection*{Le second état}

La partie A du second état (pages 254 à 262) reproduit la première, avec les
différences relevées plus haut ; la lecture n'en dit pas davantage.

\pagerange{263}{270}
\subsection*{B. Demain : les propositions}

La seconde partie « esquisse quelques-unes des propositions qui ont été
avancées, et pour la plupart discutées » — ce sont donc, pour partie au
moins, celles de la commission, qu'il expose à sa manière.

\emph{1. Deux sections.} Le département serait partagé en un \emph{Institut
de Mathématique}, qui grouperait les mathématiciens de vocation principale de
recherche et assurerait l'enseignement avancé et la formation des chercheurs,
et une \emph{École de Mathématique}, de vocation principale d'enseignement,
chargée en particulier de l'enseignement « pré-licence » aux non-mathématiciens.
Les décisions essentielles reviendraient à une \emph{Assemblée du département},
formée des maîtres de conférences et professeurs des deux sections et de
représentants des étudiants, des assistants et des maîtres-assistants. Un
membre de l'École pourrait être détaché une année sur cinq à l'Institut sans
enseignement ; un membre de l'Institut pourrait être détaché à l'École ; on
pourrait changer de section.

\emph{2. Contre la sclérose.} Année « sabbatique » pour l'École, et
obligation pour tout enseignant de renouveler son enseignement : pas plus de
trois années consécutives le même cours, sous le contrôle d'une commission
d'arbitrage. Chaque section devrait compter au moins le tiers des maîtres de
conférences et professeurs.

\emph{3. Une thèse, deux doctorats.} La thèse serait commune aux deux
orientations, et ne demanderait plus nécessairement de résultats originaux —
une mise au point originale d'une théorie avancée suffirait\footnote{Une
note du texte ajoute qu'on pourrait se borner à assouplir cette exigence
plutôt que la supprimer. Le mot « proprement » de « la thèse proprement
dite » est à demi effacé sur la copie et donné par le contexte.} ; son niveau
baisserait donc, mais elle deviendrait nécessaire et non plus suffisante.
Le doctorat aurait deux mentions. La mention « recherche » exigerait, comme
avant et même davantage, des résultats originaux d'importance indiscutable et
une personnalité capable de former des chercheurs. La mention
« enseignement » serait décernée sur critères pédagogiques : quatre à cinq ans
d'enseignement, la responsabilité complète d'un cours au moins une année, au
moins trois sujets différents, l'avis des chefs de service et celui des
étudiants recueilli chaque année, et un étudiant élu au jury.

\emph{4. Une sélection sérieuse.} Pour que la revalorisation de
l'enseignement ne soit pas « une façon pieuse de cacher une dévalorisation »
de la fonction, le doctorat d'enseignement doit être exigeant, et
l'\emph{ancienneté} ne doit jamais y être un argument. Les nominations à
l'École seraient décidées par l'Assemblée, donc avec les délégués étudiants ;
celles à l'Institut, par cooptation dans une assemblée restreinte à ses
membres.

\emph{5. L'assistant, une fonction de passage.} Dans l'idéal, être assistant
ou maître-assistant serait une étape non permanente, de quatre à six ans,
vers le doctorat et un poste de maître de conférences. On ne titulariserait
qu'après la soutenance de thèse et deux ans d'enseignement, et l'on aiderait
à se réorienter ceux qui n'y parviendraient pas. Le calcul est simple : cinq
ans de fonction d'assistant sur quarante de carrière, cela fait un huitième
d'assistants parmi les enseignants — un cinquième ou un quart pendant la
transition. « C'est à peu près la proportion inverse de celle qui est de
règle à l'heure actuelle ! » Un tel renversement suppose la fin des cours
magistraux, que les étudiants réclament et que beaucoup de professeurs ont du
mal même à imaginer.

\emph{6. Les obstacles.} Les crédits d'abord — un maître-assistant « coûte
moins cher » ; mais « les derniers évènements » permettent de juger si ces
économies sont rentables. Puis la question, plus délicate, du nombre
d'assistants actuels capables d'obtenir rapidement un doctorat
d'enseignement : beaucoup, stimulés par cette voie nouvelle, y parviendraient
sans doute ; les autres formeraient pour longtemps une « arrière-garde de
maîtres-assistants à vie ».

\emph{7. La transition et la province.} Il faudrait attribuer la majorité des
postes nouveaux à la mention « enseignement » jusqu'à atteindre un tiers au
moins. Les départements très attractifs, Paris et Orsay, devraient résister à
la tentation de « thésauriser » les meilleurs chercheurs — tentation dont il
dit qu'il serait « le dernier à sous-estimer la force » pour un mathématicien
passionné. Il propose une \emph{Commission de mathématiciens
interdépartements}, dotée d'une autorité réelle et dans certains cas d'un droit
de veto sur les nominations, pour répartir raisonnablement les chercheurs
entre les universités, à l'exclusion de toute instance ministérielle ; cela
suppose que les Parisiens renoncent à une concentration qui est un
« gaspillage par l'accumulation ».\footnote{Quelques retouches à l'encre sur
cette page (ponctuation, guillemets) sont d'une main non déterminée ; la
transcription donne le texte dans son état corrigé.}

\emph{8. Le seul point essentiel.} Jusqu'ici, l'autorité dans le département
appartenait aux seuls maîtres de conférences et professeurs, recrutés comme
chercheurs : une « aristocratie éclairée » gouvernant une armée d'assistants
quatre ou cinq fois plus nombreuse et un « sous-prolétariat » d'étudiants,
dont les besoins ont été pratiquement ignorés. Ceux qui détiennent cette
autorité y résisteront, et des discussions entre collègues, sans les
étudiants, deviendraient vite byzantines et permettraient aux immobilistes de
« noyer le poisson ». Aussi toutes les propositions qui précèdent ne sont-elles
« en dernière analyse » que des points de détail, à corriger par l'expérience
dans l'esprit de la « contestation permanente ». Un seul point de structure
compte : la \emph{cogestion} effective, c'est-à-dire la participation des
étudiants, par une représentation massive dans les instances dirigeantes, à
\emph{toutes} les décisions qui les concernent. Ce point acquis, « les
structures s'ajusteront d'elles-mêmes ; non en un jour ni en une année, mais
chaque chose en son temps ».\footnote{La page porte « contestation
permanents », faute de frappe. Pour situer : la loi d'orientation de
l'enseignement supérieur du 12 novembre 1968 (loi Edgar Faure) introduira
quelques mois plus tard la participation des étudiants aux conseils
d'université, sous une forme bien plus limitée que celle qu'il demande.}

\pagerange{272}{280}
\subsection*{Les pièces de la commission d'Orsay}

Le dossier conserve à la suite deux pièces collectives, qui ne sont pas de
lui et ne sont pas transcrites : une circulaire ronéotypée de la commission
sur les deux voies d'accès au doctorat d'État (mentions « recherche » et
« enseignement »), qui annonce une réunion « le mardi 11 juin » (pages 272
et 273) ; et le « Rapport de la Commission relative au recrutement des
enseignants » (pages 275 à 280). Selon la note de la transcription, ce
rapport partage le département en une École de Mathématique et un Laboratoire
de Recherche Mathématique sous une Assemblée du département, avec deux
mentions du doctorat, des nominations décidées par l'Assemblée ou ses
commissions, un droit de veto étudiant pour l'École, des détachements et des
changements de section : c'est, sous d'autres noms, l'architecture de la
partie B. Au pied de la circulaire, des notes au crayon, \emph{probablement}
de sa main, en reprennent les thèmes : « assistants pas à l'ancienneté »,
« nominations par étudiants », « thèse », « enseignements parallèles
(diversification) », « mandarinat », et en marge « (un peu de recherche) »,
« (décision des étudiants) ».\footnote{Plusieurs de ces mots sont de lecture
incertaine (« ancienneté », « parallèles », « diversification »), un est
illisible, et l'attribution de la main n'est pas certaine.}

\pagerange{301}{313}
\section*{Le curriculum vitae (pages 303 à 313, vers 1971)}

La page 301 est la dernière d'un résumé de cours de J.-P.~Serre pour
1990--1991, et la page 308 une lettre-type du CNRS du 4 juillet 1972 qui lui
retourne sa notice de titres et travaux ; ces deux pièces sont d'autres
mains et ne sont pas transcrites. Entre elles, et après, le dossier conserve
deux exemplaires d'un curriculum vitae dactylographié de quatre pages, non
daté. Il est postérieur à septembre 1970, puisqu'il raconte le départ de
l'IHES, et vraisemblablement antérieur à la lettre du CNRS de 1972 ; une
bande collée en tête (« Vous trouverez ci-joint l'exposé des titres travaux
des candidats à une direction de recherche ») le rattache à une candidature
au CNRS.

\subsection*{Le récit}

Il est né le 28 mars 1928 à Berlin, d'une mère allemande et d'un père
apatride émigré de Russie en 1921. Ses parents quittent l'Allemagne en 1933
et prennent part à la révolution espagnole ; il les rejoint en mai 1939. Son
père est interné en 1939, sa mère en 1940 avec lui ; son père est déporté du
camp du Vernet à Auschwitz en août 1942 et n'est pas revenu ; sa mère meurt
en 1957 d'une tuberculose contractée au camp. Après près de deux ans dans
des camps d'internement français, il est recueilli par une maison d'enfants
du Secours suisse au Chambon-sur-Lignon, où il finit ses études secondaires
en 1945.\footnote{Le texte dit « camps de concentration français » ; c'est
son terme pour les camps d'internement où il a été détenu.}

Licence à Montpellier (1945--1948), auditeur libre à l'École normale
supérieure (1948--1949), où il suit le premier séminaire Cartan sur les
faisceaux et un cours de Leray au Collège de France sur le degré topologique
de Schauder. De 1949 à 1953, à Nancy, élève de Dieudonné et de Schwartz, il
travaille sur les espaces vectoriels topologiques ; sa thèse (1953) porte sur
les produits tensoriels topologiques et les espaces nucléaires. Deux ans à
São Paulo, où il achève ces recherches et commence, sous l'influence de Serre,
à apprendre la topologie algébrique et l'algèbre homologique — « encore très
loin d'être menées à leur terme ». En 1955, à l'université du Kansas, il
unifie la théorie des foncteurs dérivés de Cartan--Eilenberg et la cohomologie
des faisceaux de Leray--Cartan (l'article du \emph{Tôhoku}), et développe une
« cohomologie non commutative » qui trouvera son cadre dans la théorie des
topos.

Chercheur au CNRS de 1950 à 1958, directeur de recherches en 1958 ;
professeur à l'IHES de 1959 à 1970. Ayant découvert à la fin de 1969 que
l'IHES recevait depuis trois ans des subventions du ministère des Armées, et
n'ayant pas réussi à entraîner ses collègues dans une action commune, il
quitte l'IHES en septembre 1970. Marié depuis 1959, père de quatre enfants,
apatride depuis 1940, il a demandé la naturalisation française au printemps
1970.

Depuis 1956, dit-il, son intérêt principal est la géométrie algébrique, et
tout le reste — topologie, géométrie analytique, algèbre homologique,
langage catégorique — lui a été subordonné, au service d'un « vaste programme
de construction » dont la première vision d'ensemble date de 1958 et qui est
poursuivi dans les \emph{Éléments} et les \emph{Séminaires}.\footnote{La page
renvoie ici à « [16, 17] » ; [17] étant un article sur les schémas abéliens,
il s'agit très probablement de [15, 16], les \emph{Éléments de géométrie
algébrique} et les \emph{Séminaires de géométrie algébrique}. Plus haut, le
renvoi « [16, SGA 4] » pour la théorie des topos est juste.} Ce programme
« est loin d'être achevé », et la fin du texte en dit la raison :

\begin{quote}
L'extraordinaire crise écologique que nous aurons à affronter dans les
décades qui viennent rend peu probable qu'il le sera jamais.
\end{quote}

Cette crise imposera des critères de valeur nouveaux, qui réduiront « à
l'insignifiance » bien des progrès scientifiques du siècle, dans la mesure où
ils restent étrangers à « l'impératif évolutionniste de notre temps : celui
de la survie ». C'est ce qui l'a conduit en juillet 1970 à cofonder le
mouvement international et interprofessionnel \emph{Survivre}. Et il
conclut que la seule valeur de son apport de mathématicien est de donner,
grâce à l'estime acquise, plus de force à son témoignage en faveur d'une
stricte subordination de toutes nos activités, scientifiques comprises, aux
impératifs de la survie et d'« un ordre stable et humain sur notre planète,
sans lequel la survie de notre espèce ne serait ni possible, ni désirable ».

Le second exemplaire (pages 310 à 313) porte, à la page 311, une correction
manuscrite de cette phrase : la crise devient « crise écologique de
civilisation dans laquelle nous sommes engagés à l'heure actuelle, et qui
devra se résoudre dans les décades qui viennent ».\footnote{« Nous aurons à
affronter » est barré d'un trait fin, « nous » biffé, une virgule ajoutée
après « viennent » ; le « que » n'est pas biffé, de sorte que la phrase
corrigée n'est pas entièrement grammaticale sur la page. « Devra » est de
lecture incertaine.} La crise n'est plus à venir : on y est.

\subsection*{La liste des publications}

Elle compte dix-neuf titres, répartis en espaces vectoriels topologiques
([1]--[7]), topologie et algèbre homologique ([8], [9]), géométrie analytique
([10], [11]) et géométrie algébrique ([12]--[19]) — dont l'article du
\emph{Tôhoku} [8], les \emph{Éléments} [15], les \emph{Séminaires} [16] et
le livre sur le complexe cotangent [19]. Le second exemplaire ajoute à la
main un vingtième titre : « Platitude d'une adhérence schématique et lemme de
Hironaka généralisé » (avec H.~Seydi), \emph{manuscripta math.} 5 (1971),
323--339.

Quelques erreurs matérielles de la liste, relevées par la transcription, sont
corrigées ici pour qui voudrait s'en servir :\footnote{Les corrections de
pagination et de titre sont de la lecture et proviennent des publications
elles-mêmes, non des pages du dossier.}
\begin{itemize}
\item [13], « Sur une note de Mattuck--Tate », \emph{J. reine angew. Math.}
  200 (1958), occupe les pages 208--215 ; la liste répète par erreur la
  pagination de [12] ;
\item dans [16], la ligne de SGA~4 manque : « Théorie des topos et
  cohomologie étale des schémas » est tapé à la suite de SGA~3, qui est
  \emph{Schémas en groupes} (avec M.~Demazure) ;
\item [9] porte « fasceaux » pour « faisceaux », et « Saõ Paulo » est écrit
  pour São Paulo.
\end{itemize}

\end{document}
